\documentclass[11pt]{amsart}
\usepackage[T1]{fontenc}
\usepackage{lmodern}
\usepackage{amsmath,amssymb,amsthm}
\usepackage[margin=1.1in]{geometry}
\usepackage{booktabs}
\usepackage{hyperref}

\newtheorem{theorem}{Theorem}[section]
\newtheorem{lemma}[theorem]{Lemma}
\newtheorem{proposition}[theorem]{Proposition}
\newtheorem{corollary}[theorem]{Corollary}
\newtheorem{conjecture}[theorem]{Conjecture}
\newtheorem{question}[theorem]{Question}
\theoremstyle{definition}
\newtheorem{definition}[theorem]{Definition}
\newtheorem{remark}[theorem]{Remark}
\newtheorem{example}[theorem]{Example}

\newcommand{\cD}{\mathcal D}
\emergencystretch=3em
\newcommand{\Good}{\mathrm{Good}}

\title[The avoidance conjecture $L_1^*$]{The avoidance conjecture $L_1^*$:\\ a local form, a
reduction to blocked instances, and two even vertices}
\author{Justin Williams}
\date{November 2026 (working draft)}

\begin{document}
\begin{abstract}
Let $G$ be a graph with at most one vertex $g$ of positive even degree. A \emph{Lov\'asz
decomposition} of $G$ is a partition of $E(G)$ into paths in which every vertex $v$ is the end of
exactly $d(v)\bmod 2$ paths; for an odd vertex $c$ let $Q_c$ be the path ending at $c$. Conjecture
$L_1^*$ asks for a Lov\'asz decomposition with $Q_c\cap A=\emptyset$ whenever $|A|\le2$ and $A$
avoids $N[c]\cup N(g)$; it implies Conjecture~EE on prescribed end edges, the open case of the
single-transition step of the star-terminal programme for Gallai's conjecture.

\begin{itemize}
\item \emph{A local form.} Computation supports a much stronger statement: some Lov\'asz
decomposition has $V(Q_c)\subseteq N[c]\cup N(g)$. Equivalently, every set $A$ disjoint from
$N[c]\cup N(g)$, of any size, can be avoided. There is no counterexample among all connected graphs
of order at most~$11$ ($119{,}593{,}670$ pairs $(G,c)$ at order $11$), nor among $45{,}714$ decided
random pairs of orders $11$--$20$. In all-odd graphs it says $V(Q_c)\subseteq N[c]$. It implies that
a graph whose only two even vertices are at distance at least three has a Lov\'asz decomposition
(also verified to order $10$, and on random graphs to order $16$).
\item \emph{A reduction.} We prove, and formalise in Lean, that a smallest counterexample to the
general form of $L_1^*$ is \emph{blocked}: $G$ has an even vertex $g$, $c\notin N(g)$, and every
neighbour of $g$ is adjacent to a vertex of $A\setminus\{g\}$. The proof is Lov\'asz's induction with
the end at $c$ protected through leaf graphs and Fan's lemma with vertex sets. Hence EE follows from
$L_1^*$ on blocked instances.
\item \emph{Further reductions} (route moves, splitting at vertices of degree two or three, a
triangle swap) restrict a smallest counterexample much further
(Corollary~\ref{cor:mincex}). At order $9$, $13.9\%$ of the instances with $|A|\le2$ are blocked,
and $0.59\%$ survive all reductions one level deep.
\end{itemize}
The reduction theorem and its corollary for EE are formalised in Lean~4 with Mathlib, with no
\texttt{sorry} and only the standard axioms. Computations are in a separate, labelled section; no
proof depends on them.
\end{abstract}
\maketitle

\section{Introduction}\label{sec:intro}

All graphs are finite and simple. A \emph{path decomposition} of $G$ partitions $E(G)$ into paths,
each with at least one edge. For a decomposition $\cD$ and a vertex $v$, $b_\cD(v)$ is the number of
paths of $\cD$ with an end at $v$; always $b_\cD(v)\equiv d(v)\pmod 2$. A vertex is \emph{even} if
its degree is positive and even, and \emph{odd} if its degree is odd. Following~\cite{RD}, $\cD$ is
a \emph{Lov\'asz decomposition} if $b_\cD(v)=d(v)\bmod2$ for every $v$. Lov\'asz's
theorem~\cite{L68} says that a graph with at most one even vertex has one. For an odd vertex $c$,
$Q_c$ is the path of $\cD$ ending at $c$. Throughout, $g$ denotes the even vertex when there is one,
and $N(g)=\emptyset$ otherwise.

\begin{definition}
An \emph{instance} is a triple $(G,c,A)$: $G$ has at most one even vertex, $c$ is odd, and
$A\subseteq V(G)$ is \emph{admissible}: $A\cap(N[c]\cup N(g))=\emptyset$. ($g\in A$ is allowed when
$g\notin N(c)$.) The instance is \emph{good} if some Lov\'asz decomposition of $G$ has
$V(Q_c)\cap A=\emptyset$.
\end{definition}

\begin{conjecture}[$L_1^*$~{\cite[Conjecture~7.3]{PE}}]\label{conj:L1}
Every instance with $|A|\le2$ is good.
\end{conjecture}

\begin{conjecture}[$L_1^*$, general form]\label{conj:gen}
Every instance is good.
\end{conjecture}

\begin{conjecture}[local form]\label{conj:loc}
If $G$ has at most one even vertex and $c$ is odd, some Lov\'asz decomposition has
$V(Q_c)\subseteq N[c]\cup N(g)$.
\end{conjecture}

Conjectures~\ref{conj:gen} and~\ref{conj:loc} are equivalent. The set
$A_{\max}=V\setminus(N[c]\cup N(g))$ is admissible and contains every admissible set, and avoiding
$A_{\max}$ means $V(Q_c)\subseteq N[c]\cup N(g)$. (If $g\notin N(c)$ then $g\in A_{\max}$, so $Q_c$
avoids $g$.)

\emph{Why it matters.} Conjecture~EE~\cite[Conjecture~1.5]{PE} says: if $z$ is the only even vertex
and $u\in N(z)$, then $uz$ is the last edge of $Q_u$ in some Lov\'asz decomposition. It is proved
when $u$ and $z$ have at most one common neighbour~\cite[Theorem~5.1]{PE}, and in general it follows
from Conjecture~\ref{conj:L1}~\cite[Corollary~7.4]{PE} (Lean: \texttt{ee\_of\_L1Star}). EE is the
single-transition case of the star-terminal programme, which reduces Gallai's conjecture to
statements about Lov\'asz decompositions (Lean: \texttt{gallai\_of\_stl}).

\subsection*{Results}
\begin{enumerate}
\item The local form (Conjecture~\ref{conj:loc}) has no counterexample in our data
(Section~\ref{sec:comp}). Its consequences (Section~\ref{sec:local}): in all-odd graphs,
$V(Q_c)\subseteq N[c]$; and the two-even conjecture (Conjecture~\ref{conj:te3}), which says that two even
vertices at distance at least three do not prevent a Lov\'asz decomposition. The hypothesis
$A\cap N(g)=\emptyset$ cannot be weakened to ``$A$ separates $c$ from $g$'' (Example~\ref{ex:sep}),
and the local form fails with two even vertices even where Lov\'asz decompositions exist
(Remark~\ref{rem:pe}).
\item Theorem~\ref{thm:blocked}: a smallest counterexample to Conjecture~\ref{conj:gen} is blocked
(formal). Corollary~\ref{cor:eeblocked}: EE follows from the blocked case (formal).
\item Section~\ref{sec:more}: five more reductions and the resulting structure of a smallest
counterexample (Corollary~\ref{cor:mincex}).
\item Section~\ref{sec:comp}: an exact oracle, validated against brute force; exhaustive results to
order~$11$; random graphs to order $16$; a census of blocked and residual instances.
\end{enumerate}

\section{Preliminaries}\label{sec:prelim}

For a decomposition $\cD$ and a vertex $w$ let $Z_\cD(w)=\{v\in N(w):b_\cD(v)=0\}$.

\begin{lemma}[{Fan's lemma with vertex sets~\cite{Fan05},~\cite[Lemma~2.2]{RD}}]\label{lem:fan}
Let $\cD$ be a path decomposition of $K$, let $t\ne w$ be non-adjacent, and suppose
$b_\cD(t)>|Z_\cD(w)|$. Then $K+tw$ has a path decomposition $\cD^+$ with
$b_{\cD^+}=b_\cD-\mathbf 1_t+\mathbf 1_w$ in which every path has its vertex set inside $V(P)\cup\{w\}$
for some path $P$ of $\cD$.
\end{lemma}

\emph{Leaf graphs.} For an edge $xy$ let $H(K;x,y)$ be $K-xy$ with a new vertex $\ell$ and the edge
$y\ell$; the leaf $\ell$ stands for $x$. For a vertex $x$ let $K^{+x}$ be $K$ with a new vertex
$\ell$ and the edge $x\ell$ (a leaf hung at $x$).

\begin{lemma}[leaf merge~{\cite[Lemma~12]{MS}}, with avoidance]\label{lem:merge}
Let $\cD_H$ be a path decomposition of $H=H(K;x,y)$ in which no path contains $\ell$ together with
$x$, and no path contains $\ell$ together with a vertex of a set $A\not\ni x$. Suppose $b_{\cD_H}(x)=0$
and $b_{\cD_H}(\ell)=1$. Identifying $\ell$ with $x$ gives a decomposition $\cD$ of $K$ with
$b_\cD(x)=1$, $b_\cD=b_{\cD_H}$ elsewhere, and $V(Q_x)\cap A=\emptyset$.
\end{lemma}

\begin{proof}
The merge is~\cite[Lemma~12]{MS}; it maps paths to paths and the path through $y\ell$ to a path
through $yx$ with $\ell$ replaced by $x$. That path ends at $x$ (as $\ell$ is a leaf), and it is the
only one, since $b_\cD(x)=0+1$. Its vertices other than $x$ are vertices of the leaf path, which
avoids $A$.
\end{proof}

\begin{lemma}[a leaf hung at $c$]\label{lem:hang}
\emph{(i)} If $\cD$ decomposes $K$ and $b_\cD(c)=1$, extending $Q_c$ by $c\ell$ gives a
decomposition of $K^{+c}$ with $b(c)=0$, $b(\ell)=1$, the other end counts of $\cD$, and the vertex
sets of $\cD$ except that $Q_c$ gains $\ell$. \emph{(ii)} Conversely, let $\cD^+$ decompose $K^{+c}$
with $b(c)=0$, $b(\ell)=1$, and let no path contain $\ell$ together with a vertex of $A$. Deleting
$\ell$ gives a decomposition $\cD$ of $K$ with $b_\cD(c)=1$, $b_\cD=b_{\cD^+}$ elsewhere, and
$V(Q_c)\cap A=\emptyset$.
\end{lemma}

\begin{proof}
(i) is clear. (ii) Deleting the leaf shortens the path $\ell\,c\cdots$ to a path starting at $c$
(it is not the single edge $\ell c$, since $c$ has no end in $\cD^+$). This path avoids $A$ and ends
at $c$; every other path is unchanged and has no end at $c$.
\end{proof}

\begin{lemma}[protected transfer]\label{lem:protect}
Let $\cD$ decompose $K$, let $t\ne w$ be non-adjacent with $b_\cD(t)=1$ and $Z_\cD(w)=\emptyset$, and
let $c\notin N(w)\cup\{t,w\}$ have $b_\cD(c)=1$ with $V(Q_c)\cap A=\emptyset$, where $w\notin A$. Then
$K+tw$ has a decomposition $\cD^+$ with $b_{\cD^+}=b_\cD-\mathbf 1_t+\mathbf 1_w$ and
$V(Q_c)\cap A=\emptyset$.
\end{lemma}

\begin{proof}
Hang a leaf at $c$ (Lemma~\ref{lem:hang}(i)). In $K^{+c}$ the vertex $c$ has no end but is not a
neighbour of $w$, and $\ell$ is not a neighbour of $w$; so $Z(w)$ is unchanged and empty.
Lemma~\ref{lem:fan} in $K^{+c}$ restores $tw$; every new path lies in an old one together with $w$, so
still no path contains $\ell$ and a vertex of $A$. Delete $\ell$ (Lemma~\ref{lem:hang}(ii)).
\end{proof}

This is a vertex-set form of the protected-end lemma of~\cite[Lemma~4.1]{PE}.

\section{The local form and its consequences}\label{sec:local}

\begin{proposition}\label{prop:allodd}
In a graph with no even vertex, Conjecture~\ref{conj:loc} says $V(Q_c)\subseteq N[c]$. If $d(c)=1$
it holds: $Q_c$ can be the edge at $c$.
\end{proposition}

\begin{proof}
If $N(c)=\{y\}$, then $G-cy$ has at most one even vertex ($y$), so it has a Lov\'asz decomposition;
add the path $cy$.
\end{proof}

\begin{conjecture}[two even vertices]\label{conj:te3}
If the only two even vertices $x,y$ of $G$ are at distance at least three, $G$ has a decomposition
into paths with $b(v)=d(v)\bmod2$ for every $v$.
\end{conjecture}

\begin{proposition}\label{prop:te3}
Conjecture~\ref{conj:loc} implies Conjecture~\ref{conj:te3}.
\end{proposition}

\begin{proof}
Hang a leaf $w$ at $y$. In $G^{+y}$ the vertex $y$ is odd, so $x$ is the only even vertex. Apply
Conjecture~\ref{conj:loc} to $c=w$: some Lov\'asz decomposition has
$V(Q_w)\subseteq\{w,y\}\cup N(x)$. The second vertex of $Q_w$ is $y$. A third vertex would lie
in $N(y)\cap N(x)$, which is empty. So $Q_w=wy$, and as $b(y)=1$, $y$ ends no other path.
Delete $Q_w$: what remains decomposes $G$ with $b(y)=0$, $b(x)=0$, and $b(v)=d(v)\bmod2$ elsewhere.
\end{proof}

Conjecture~\ref{conj:te3} holds for all connected graphs of order at most $10$
(Section~\ref{sec:comp}). It is false at distance $1$ and $2$: $22{,}100$ of the $562{,}856$ and
$21{,}258$ of the $501{,}456$ connected graphs of order $10$ with two even vertices at distance $1$,
respectively $2$, have no such decomposition. The more general statement that \emph{packed} even vertices (pairwise at distance at
least three) allow a Lov\'asz decomposition also holds to order $10$.

\begin{remark}[two even vertices and the local form]\label{rem:pe}
The local form does not extend to two even vertices. Among the connected graphs of order $10$ with two
even vertices $g_1,g_2$ at distance at least three, $61$ of the $404{,}488$ pairs $(G,c)$ have no
decomposition with $V(Q_c)\subseteq N[c]\cup N(g_1)\cup N(g_2)$, though each graph has a Lov\'asz
decomposition. One example is \texttt{GCQbVg} with $c=1$ at order $8$.
\end{remark}

\begin{remark}[one-sided separation]\label{rem:ball}
Route descent~\cite[Conjecture~8.2]{RD} conjectures a different region. With at most one even vertex,
some Lov\'asz decomposition should have every vertex of $Q_c$ at distance at most two from $c$:
$V(Q_c)\subseteq B_2(c)$. For a leaf $c$ with neighbour $y$ this reads $V(Q_c)\subseteq\{c\}\cup N[y]$,
and through the leaf trick it gives $V(Q_c)\subseteq N[c]$ in all-odd graphs, as
Proposition~\ref{prop:allodd} does. That region does not depend on $g$, which is why Lov\'asz's
induction nearly proves the leaf case there. The one residual configuration has the even vertex at
distance three from $c$ and all its neighbours at distance two. It is the same tight case of Fan's
inequality as the blocked case below. The two regions can be combined. Some Lov\'asz decomposition has
\[V(Q_c)\subseteq N[c]\cup\bigl(N(g)\cap B_2(c)\bigr),\]
which implies Conjecture~\ref{conj:loc} and~\cite[Conjecture~8.2]{RD}, for every connected graph of
order at most $9$ (\texttt{ball\_test.py}, \texttt{out/ball\_5\_9.txt}).
\end{remark}

\begin{example}[the hypothesis $A\cap N(g)=\emptyset$]\label{ex:sep}
In a star $K_{1,4}$ with centre $g$ and $c$ a leaf, $Q_c=c\,g\,a$ for another leaf $a$, so $A$ cannot
contain the three other leaves. The condition cannot be replaced by ``$A$ separates $c$ from $g$''.
Take vertices $0,\dots,4$ with edges $02,03,04,14,23,24$ (graph6 \texttt{DTw}). Here $g=3$, $c=1$ and
$A=\{0,2\}=N(g)$, so $A$ separates $c$ from $g$. The path $Q_1$ must be $14$, and then $G-14$ must be a
single path from $0$ to $2$ with five edges on four vertices, which is impossible. At order~$9$,
$3{,}162$ of the $47{,}526$ pairs $(c,A)$ with $|A|\le4$, $A\cap N[c]=\emptyset$, $A$ separating $c$
from $g$ and $A\cap N(g)\neq\emptyset$ fail.
\end{example}

\section{Reduction to blocked instances}\label{sec:blocked}

\begin{definition}
An instance $(G,c,A)$ is \emph{blocked} if $G$ has an even vertex $g$, $c\notin N(g)$, and every
neighbour of $g$ has a neighbour in $A\setminus\{g\}$.
\end{definition}

We use Lov\'asz's measure $\nu(G)=\sum_vd(v)+[G\text{ has no even vertex}]$.

\begin{theorem}\label{thm:reductions}
Let $(G,c,A)$ be an instance.
\begin{enumerate}
\item[(a)] If $c\in N(g)$, then $(G-gc,\,g,\,A\cup\{c\})$ is an instance, and if it is good so is
$(G,c,A)$.
\item[(b)] If $c\notin N(g)$ and $u\in N(g)\setminus N(c)$ has no neighbour in $A\setminus\{g\}$, then $(G-gu,\,c,\,A)$
is an instance, and if it is good so is $(G,c,A)$.
\item[(c)] If $u\in N(g)\cap N(c)$ has no neighbour in $A\setminus\{g\}$, then
$(G-cu-gu,\,u,\,A\cup\{c\})$ is an instance, and if it is good so is $(G,c,A)$.
\item[(d)] If $G$ has no even vertex and $y\in N(c)$, $d(c)\ge3$, then $(H(G;c,y),\,\ell,\,A\cup\{c\})$
is an instance, and if it is good so is $(G,c,A)$.
\end{enumerate}
In (a)--(c) the new graph has smaller $\nu$; in (d), $\nu(H)=\nu(G)-1$.
\end{theorem}

\begin{proof}
Write $\deg'$ for degrees in the new graph. In each case only the stated vertices change degree.

(a) In $G'=G-gc$: $g$ is odd, $c$ is even (or isolated), so $c$ is the only possible even vertex,
and $g$ is the new terminal. $A\cup\{c\}$ is admissible: $c\notin N'[g]$ as $gc$ was deleted;
$g\notin A$ as $g\in N(c)$; $A\cap N(g)=\emptyset$; $N'(c)\subseteq N(c)$ misses $A$, and
$c\notin N(c)$. Let $\cD'$ be good: $Q'_g$ avoids $A\cup\{c\}$. Hang the leaf at $g$ by extending
$Q'_g$; the result decomposes $H(G;c,g)$, where $\ell$ stands for $c$, and has $b(c)=0$, $b(\ell)=1$.
No path contains $\ell$ with $c$ or with a vertex of $A$. Lemma~\ref{lem:merge} gives a decomposition
of $G$ with $b(c)=1$, $b(g)=b'(g)-1=0$, the Lov\'asz counts elsewhere, and $Q_c$ disjoint from $A$.
This is ``append $gc$ to $Q'_g$''.

(b) In $G'=G-gu$, $u$ is the only possible even vertex, $c$ is odd (it is neither $g$ nor $u$, since
$c\notin N(g)$), and $A$ is admissible: a vertex of
$A$ in $N'(u)$ would be a neighbour of $u$ in $A\setminus\{g\}$. Given a good $\cD'$, apply
Lemma~\ref{lem:protect} with $t=g$, $w=u$: $b'(g)=1$; every neighbour of $u$ in $G'$ is odd, so
$Z(u)=\emptyset$; $c\notin N(u)$; and $u\notin A$ since $u\in N(g)$.

(c) In $G''=G-cu-gu$: $u$ has lost two edges and is odd, $c$ and $g$ have lost one each, so $c$ is
the only possible even vertex. $A\cup\{c\}$ is admissible for $u$: $c\notin N''[u]$;
$A\cap N''[u]=\emptyset$ since $u\notin A$ and $u$ has no neighbour in $A\setminus\{g\}$, while
$gu$ was deleted; $A\cap N''(c)=\emptyset$. Let $\cD''$ be good. Hang a leaf at $u$ by extending
$Q''_u$. This decomposes $H(G-gu;c,u)$, where $\ell$ stands for $c$; no path contains $\ell$ with $c$
or with a vertex of $A$; $b(u)=0$, $b(g)=1$. In it, $u$ and $g$ are non-adjacent and every neighbour
of $u$ has an end ($\ell$, and odd vertices; $c$ is no longer a neighbour). Lemma~\ref{lem:fan} with
$t=g$, $w=u$ gives a decomposition of $H(G;c,u)$. It keeps both separations, because $u\notin
A\cup\{c,\ell\}$, and has $b(u)=1$, $b(g)=0$. Merge (Lemma~\ref{lem:merge}).

(d) In $H=H(G;c,y)$, $c$ has degree $d(c)-1\ge2$ and is the only even vertex, and $\ell$ is a
leaf. $A\cup\{c\}$ is admissible for $\ell$: $N_H[\ell]=\{\ell,y\}$ misses it because $y\in N(c)$;
$N_H(c)\subseteq N(c)$ misses $A$. A good decomposition of $H$ has $Q_\ell$ disjoint from $A\cup\{c\}$,
and since $\ell$ is a leaf, every path through $\ell$ is $Q_\ell$. Merge (Lemma~\ref{lem:merge}).

Measures. The sum of degrees drops by $2$ in (a) and (b) and by $4$ in (c); the indicator rises by at
most one. In (d) the sum is unchanged and $H$ has an even vertex while $G$ has none.
\end{proof}

\begin{theorem}\label{thm:blocked}
If every blocked instance is good, every instance is good. Equivalently, a counterexample to
Conjecture~\ref{conj:gen} with $\nu(G)$ smallest is blocked.
\end{theorem}

\begin{proof}
Induction on $\nu(G)$. If $G$ has an even vertex $g$, either $c\in N(g)$ (Theorem~\ref{thm:reductions}(a)),
or some $u\in N(g)$ has no neighbour in $A\setminus\{g\}$ ((b) or (c)), or the instance is blocked.
If $G$ has no even vertex, $d(c)=1$ (Proposition~\ref{prop:allodd}) or $d(c)\ge3$ ((d)).
\end{proof}

\begin{corollary}\label{cor:eeblocked}
If every blocked instance is good, Conjecture~EE holds.
\emph{(Lean: \texttt{ee\_of\_\allowbreak l1Blocked}.)}
\end{corollary}

\begin{proof}
Theorem~\ref{thm:blocked}, then~\cite[Corollary~7.4]{PE}.
\end{proof}

\begin{remark}
The instances that the proof of EE needs are $(K-\{uz,uw,wz\},\,w,\,\{u,z\})$ for a common
neighbour $w$ of $u$ and $z$. If $u$ has a neighbour outside $N[z]$, one may use
$(K-uz-uw,\,z,\,\{u\})$ instead~\cite{PE}. The first instance is never blocked when $N(u)\setminus\{z\}\subseteq N(z)$.
Being blocked would mean $N(z)\setminus\{u\}\subseteq N(u)$; then $u$ and $z$ would have the same
neighbours apart from each other, hence the same degree, but $d(u)$ is odd and $d(z)$ even. Blocked
instances do arise one or two reductions later. In the double $K_4$ (\texttt{FQhVw}, two $K_4$'s
sharing $z$), every branch of the reductions of Theorem~\ref{thm:reductions} alone, from the
instances $(K-uz,z,\{u\})$ and from the triangle instances, ends in a blocked instance with
$d(g)=2$. No reduction of Section~\ref{sec:more} applies to any of these blocked instances (audit
A23). One of them is $K-\{02,06,26,16\}$ with $g=1$, $c=2$, $A=\{0,6\}$. It is good, with $Q_c=24$.
Interleaving the reductions of Section~\ref{sec:more} earlier in the branches certifies the six
instances $(K-uz,z,\{u\})$ (\texttt{cert\_ee.py}), but not the twelve triangle instances. With the
same interleaving, EE is certified at every pair of order at most $7$ except those of $K_7$ minus a
perfect matching of $K_6$, that is, the cocktail party graph $\mathrm{CP}(3)$ with an apex $z$
(Section~\ref{sec:comp}).
\end{remark}

\section{Further reductions}\label{sec:more}

The following reductions are proved here, not formalised. Each maps an instance to an instance that
is smaller in $(\nu(G),|A|)$ lexicographically (for Lemma~\ref{lem:boundary}, when
$\partial C\neq A$), and transfers goodness back.

\begin{lemma}[boundary]\label{lem:boundary}
Let $C$ be the component of $G-A$ containing $c$ and $\partial C=N(C)\subseteq A$. Then
$(G,c,\partial C)$ is an instance, and if it is good so is $(G,c,A)$.
\end{lemma}

\begin{proof}
$\partial C\subseteq A$ is admissible. A path from $c$ that avoids $\partial C$ stays in $C$, which
misses $A$.
\end{proof}

\begin{proposition}[route to $c$]\label{prop:route}
Let $G$ have an even vertex $g\notin A$, let $c\notin N(g)$, and let $R=g\,r_1\cdots r_{k-1}\,c$ be a
shortest $g$--$c$ path in $G-A$ (so $k\ge2$).
Then $(G-E(R),\,g,\,A\cup\{r_1,\dots,r_{k-1},c\})$ is an instance, and if it is good so is
$(G,c,A)$.
\end{proposition}

\begin{proof}
In $G'=G-E(R)$, $g$ and $c$ change parity and the interior of $R$ does not, so $c$ is the only
possible even vertex. Admissibility: $R$ has no chords, being shortest in $G-A$, so
$N'[g]\cap V(R)=\{g\}$ and $N'(c)\cap V(R)=\emptyset$. Also $A\cap N'[g]=\emptyset$ ($g\notin A$,
$A\cap N(g)=\emptyset$) and $A\cap N'(c)=\emptyset$. If $Q'_g$ avoids $A\cup(V(R)\setminus\{g\})$, then
$Q'_g$ and $R$ share only $g$, an end of both, so $Q'_g\cup R$ is a path. Replacing $Q'_g$ by it gives a Lov\'asz decomposition of $G$: $g$ loses its
end and $c$ gains one. Its path ending at $c$ is $Q'_g\cup R$, which misses $A$.
\end{proof}

\begin{proposition}[Fan route]\label{prop:fanroute}
Let $G$ have an even vertex $g$, and let $R=g\,e_1\cdots e_k$ ($k\ge1$) be a path with $c\notin V(R)$,
$e_i\notin A$ for $i=1,\dots,k$, and $N(e_k)\cap A\subseteq\{e_{k-1}\}$ (where $e_0=g$). Suppose
$e_i\notin N(c)$ for $i=1,\dots,k$, or $d(c)=1$.
Then $(G-E(R),\,c,\,A)$ is an instance, and if it is good so is $(G,c,A)$.
\end{proposition}

\begin{proof}
Let $K_i=G-\{e_0e_1,\dots,e_{i-1}e_i\}$. In $K_i$, $e_i$ is the only possible even vertex, since the
other route vertices lose $0$ or $2$ edges and $g$ loses one. Admissibility in $K_k$:
$A\cap N_{K_k}(e_k)\subseteq A\cap(N(e_k)\setminus\{e_{k-1}\})=\emptyset$, and $N_{K_k}[c]=N[c]$.
Starting from a good decomposition of $K_k$, for $i=k,\dots,1$ restore $e_{i-1}e_i$ by Fan's lemma
with $t=e_{i-1}$ ($b=1$) and target $e_i$; $Z(e_i)=\emptyset$ since every other vertex of $K_i$ is
odd. If $e_i\notin N(c)$, Lemma~\ref{lem:protect} keeps $Q_c$ off $A$, as $e_i\notin A$. If $d(c)=1$
and $e_i$ is the neighbour of $c$, argue in the leaf graph, as in the proof of~\cite[Lemma~4.1]{PE}.
(1) In $H=H(K_i;c,e_i)$ the vertex $c$ is isolated, and replacing $Q_c=c\,e_i\cdots$ by
$\ell\,e_i\cdots$ gives a decomposition of $H$ in which $\ell$ ends the leaf path, which avoids $A$.
(2) The neighbours of $e_i$ in $H$ are $\ell$, which has an end, and odd vertices of $K_i$, so
$Z_H(e_i)=\emptyset<1=b(e_{i-1})$. (3) Lemma~\ref{lem:fan} in $H$ restores $e_{i-1}e_i$ and gives a
decomposition of $H(K_{i-1};c,e_i)$. Every path lies in an old one together with $e_i$, so the leaf
path still avoids $A$ (it already contains $e_i$). (4) Lemma~\ref{lem:merge} merges $\ell$ back into
$c$, which lies on no path of $H$.
\end{proof}

Theorem~\ref{thm:reductions}(b),(c) are the cases $k=1$ ((c) in a stronger form that also works
when $d(c)\ge3$).

\begin{proposition}[degree two at $g$]\label{prop:deg2}
Let $N(g)=\{u,v\}$ and $c\notin N(g)$.
\begin{enumerate}
\item[(i)] If $uv\notin E$, let $G'=G-gu-gv+uv$, in which $g$ is isolated and no vertex is even. If
$g\notin A$, $(G',c,A)$ is an instance; if $g\in A$ and $u\notin N(c)$, $(G',c,(A\setminus\{g\})\cup\{u\})$
is an instance. In either case, if it is good so is $(G,c,A)$.
\item[(ii)] If $uv\in E$ and $g\in A$, then $(G,c,A\setminus\{g\})$ is an instance, and if it is good
so is $(G,c,A)$.
\end{enumerate}
\end{proposition}

\begin{proof}
(i) Admissibility needs only $N'(c)=N(c)$, since $G'$ has no even vertex. In a good decomposition of
$G'$, replace the edge $uv$ by the path $u\,g\,v$; $g$ lies on no other path. The ends are unchanged.
If the path through $uv$ was $Q'_c$, it now contains $g$. That is allowed when $g\notin A$, and it
cannot happen when $u\in A'$.

(ii) Take a Lov\'asz decomposition with $Q_c$ disjoint from $A\setminus\{g\}$. If $g\in V(Q_c)$, then
$g$ is interior, so $Q_c=X\,u\,g\,v\,Y$, and the edge $uv$ lies on another path $P$, since $u,v$ occur
once on $Q_c$ and not consecutively. $g\notin V(P)$. Exchange: $Q_c\mapsto X\,u\,v\,Y$ and
$P\mapsto P$ with $uv$ replaced by $u\,g\,v$. The ends are unchanged and $Q_c$ loses $g$.
\end{proof}

\begin{proposition}[degree three at $c$]\label{prop:deg3}
Let $N(c)=\{y_0,y_1,y_2\}$ with $c\notin N(g)$, $y_1y_2\notin E$ and $y_1\notin N(g)$. Let
$G'=G-cy_1-cy_2+y_1y_2$. Then $(G',c,A\cup\{y_1\})$ is an instance, and if it is good so is $(G,c,A)$.
\end{proposition}

\begin{proof}
Degrees, hence parities, are unchanged except that $d'(c)=1$. Admissibility: $y_1\notin N'[c]=\{c,y_0\}$
and $y_1\notin N'(g)=N(g)$. In a good decomposition of $G'$, $c$ is a leaf, so the only path through
$c$ is $Q'_c$, which avoids $y_1$. Hence the path through $y_1y_2$ avoids $c$; replace $y_1y_2$ by
$y_1\,c\,y_2$.
\end{proof}

\begin{corollary}[a smallest counterexample]\label{cor:mincex}
Let $(G,c,A)$ be a counterexample to Conjecture~\ref{conj:gen} with $\nu(G)$ smallest, and then $|A|$
smallest. Then:
\begin{enumerate}
\item $G$ has an even vertex $g$, $c\notin N(g)$, and every neighbour of $g$ has a neighbour in
$A\setminus\{g\}$ (Theorem~\ref{thm:blocked});
\item $A=\partial C$ for the component $C\ni c$ of $G-A$ (Lemma~\ref{lem:boundary});
\item if $g\notin A$, then $A$ separates $c$ from $g$ (Proposition~\ref{prop:route});
\item every path $g\,e_1\cdots e_k$ ($k\ge1$) whose vertices after $g$ avoid $A\cup\{c\}$ (and
also avoid $N(c)$, unless $d(c)=1$) ends at a vertex $e_k$ with a neighbour in $A\setminus\{e_{k-1}\}$
(Proposition~\ref{prop:fanroute});
\item if $d(g)=2$, $N(g)=\{u,v\}$, then either $uv\in E$, $g\notin A$, and $A$ separates $c$ from
$g$; or $uv\notin E$, $g\in A$, and $u,v\in N(c)$ (Proposition~\ref{prop:deg2} and item 3);
\item if $d(c)=3$, every two non-adjacent neighbours of $c$ lie in $N(g)$
(Proposition~\ref{prop:deg3});
\item $A\cup\{c\}$ is not packed (pairwise distance at least three), by the separation
theorem~\cite[Theorem~6.1]{RD}.
\end{enumerate}
\end{corollary}

\section{Computations}\label{sec:comp}

\emph{This section is computational. No proof above depends on it.} The scripts are in
\nolinkurl{manuscripts/l1star/scripts/}, and their outputs in \nolinkurl{scripts/out/}.

\subsection*{Oracle and validation}
\texttt{l1s.c} decides whether an instance is good. Its search core is the oracle of~\cite{PE}
(\texttt{pee.c}, validated there against brute force and an independent oracle). $Q_c$ is grown first
from $c$, never entering $A$. The remaining edges are then decomposed exactly with prescribed end
counts, using a failure memo and parity, degree and component pruning. \texttt{xcheck\_brute.py}
compares it with an independent brute force (\texttt{brute.py}, which enumerates all transition
systems) on every connected graph of order $4$--$7$ with at most one even vertex, every odd $c$
and every subset $A\subseteq V\setminus\{c\}$, admissible or not. There were no mismatches: $64$,
$448$, $2{,}496$ and $49{,}920$ queries at orders $4$, $5$, $6$, $7$ (\texttt{out/xcheck\_4\_6.txt},
\texttt{out/xcheck\_7.txt}).

\subsection*{The local form, exhaustively}
Mode \texttt{max} tests $A_{\max}=V\setminus(N[c]\cup N(g))$ (vertices of positive degree) for every
connected graph with at most one even vertex and every odd $c$. A pass implies every admissible $A$.
For connected graphs, odd order means exactly one even vertex and even order means none.

\begin{center}
\begin{tabular}{rrrr}
\toprule
order & graphs & pairs $(G,c)$ & failures\\
\midrule
$5$ & $7$ & $28$ & $0$\\
$6$ & $13$ & $78$ & $0$\\
$7$ & $130$ & $780$ & $0$\\
$8$ & $224$ & $1{,}792$ & $0$\\
$9$ & $11{,}936$ & $95{,}488$ & $0$\\
$10$ & $32{,}851$ & $328{,}510$ & $0$\\
$11$ & $11{,}959{,}367$ & $119{,}593{,}670$ & $0$\\
\bottomrule
\end{tabular}
\end{center}
Order $11$ ran as three \texttt{geng} slices (\texttt{run\_n11.sh}, \texttt{out/max11\_*.txt}); the
graph count equals that of the order-$11$ EE run of~\cite{PE}. In $4{,}914{,}097$ of the pairs at
order $11$, $A_{\max}$ is empty.

\subsection*{Random graphs}
\texttt{rand\_max.py 3000 7 5000000 11 12 13 14 15 16}: $G(n,p)$ with $p\in[0.15,0.6]$, parities
repaired by toggling edges between pairs of even vertices, connected; $3{,}000$ graphs ($483$--$521$
per order), $38{,}896$ pairs $(G,c)$: $38{,}697$ decided, $0$ failures, $199$ undecided (node
budget $5\cdot10^6$). Output: \texttt{out/rand\_max\_11\_16.txt}. With
\texttt{PLO=0.12 PHI=0.35 rand\_max.py 400 21 2000000 17 18 19 20}: $400$ graphs, $7{,}202$ pairs,
$7{,}017$ decided, $0$ failures, $185$ undecided (\texttt{out/rand\_max\_17\_20.txt}). For
Conjecture~\ref{conj:te3}, \texttt{rand\_te3.py 600 11 3000000 2 12 13 14 15 16} gives $360$ graphs
(two even vertices at distance at least three, even orders $12$--$16$): $358$ decomposable, $2$
undecided. \texttt{rand\_te3.py 400 12 3000000 3 13 15} gives $320$ graphs with three packed even
vertices: $318$ decomposable, $2$ undecided (\texttt{out/rand\_te3\_k*.txt}).

\subsection*{Necessity and shape}
\begin{itemize}
\item Dropping $A\cap N(g)=\emptyset$ (mode \texttt{max~3~n}) gives failures at orders $5$ and $7$,
all minimal failing sets meeting $N(g)$.
\item Mode \texttt{gen~4} at order $9$ tests all $A\subseteq V\setminus N[c]$ with $|A|\le4$:
$350{,}934+213{,}062+1{,}750$ admissible instances, $0$ failures; Example~\ref{ex:sep} for the
separated non-admissible ones.
\item The shortest possible $Q_c$ is often long: at order $8$ some pairs need a Hamiltonian $Q_c$, as
in $K_8$. The vertices of $Q_c$ outside $N[c]$ (inside $N(g)$) can number five at order $7$.
\end{itemize}

\subsection*{Two even vertices}
Mode \texttt{two}: connected graphs with exactly two even vertices; does a decomposition with
$b=d\bmod2$ exist?
\begin{center}
\begin{tabular}{rrrr}
\toprule
order & distance $1$ & distance $2$ & distance $\ge3$\\
\midrule
$6$ & $17/26$ & $10/18$ & $3/3$\\
$8$ & $1{,}191/1{,}386$ & $932/1{,}109$ & $173/173$\\
$10$ & $540{,}756/562{,}856$ & $480{,}198/501{,}456$ & $50{,}561/50{,}561$\\
\bottomrule
\end{tabular}
\end{center}
Packed even vertices (mode \texttt{packedev}): every graph is decomposable (orders $6$, $8$, $9$,
$10$; order $7$ has no graph with two packed even vertices; at order $9$, $52$ graphs with three
packed even vertices). Mode \texttt{peloc}: Remark~\ref{rem:pe}.

\subsection*{Census of reductions}
\texttt{blocked\_count.py~2} classifies the instances with $1\le|A|\le2$ (connected graphs with
one even vertex; $A$ of vertices of positive degree) by the case of Theorem~\ref{thm:reductions}
that the proof of Theorem~\ref{thm:blocked} uses. The priority is (a), then (c), then (b), so row (b)
counts the instances with a free $u$ but none in $N(c)$:
\begin{center}
\begin{tabular}{lrr}
\toprule
& order $7$ & order $9$\\
\midrule
instances & $1{,}868$ & $411{,}948$\\
(a) $c\in N(g)$ & $635$ & $144{,}284$\\
(b) free $u\notin N(c)$ & $347$ & $68{,}678$\\
(c) free $u\in N(c)$ & $618$ & $141{,}772$\\
blocked, $g\in A$ & $86$ & $13{,}672$\\
blocked, $g\notin A$ & $182$ & $43{,}542$\\
\bottomrule
\end{tabular}
\end{center}
\texttt{irred\_all.py~2} applies all reductions of Sections~\ref{sec:blocked}--\ref{sec:more} one
level deep to the $267{,}664$ instances with $c\notin N(g)$. At order $9$, $2{,}414$ instances
remain: $0.59\%$ of the $411{,}948$ instances of the census, and $0.90\%$ of those with $c\notin N(g)$.
All of them have $d(g)\in\{2,4\}$;
$2{,}219$ have $g\in A$, and the others have $A$ separating $c$ from $g$. All are good, like every
instance.

\texttt{cert.py} searches recursively for a certificate, trying shortest routes and a bounded
number of choices. Its reductions are Theorem~\ref{thm:reductions}, Lemma~\ref{lem:boundary} and
Propositions~\ref{prop:route}--\ref{prop:deg3}. It also uses one composite move, a Fan route from $g$
to a neighbour $y$ of $c$ with $N[y]\cap A=\emptyset$, followed by Theorem~\ref{thm:reductions}(c);
the counts below are the same without it (audit A23). Its base cases are: $A$ misses the component of
$c$; $A\cup\{c\}$ is packed (the separation theorem~\cite[Theorem~6.1]{RD}); no even vertex and $A$
packed~\cite[Theorem~6.3]{RD}; no even vertex and $d(c)=1$. \texttt{cert\_ee.py} does the same from
the EE instances $(K-uz,z,\{u\})$. At order $6$ every instance with $|A|\le2$ is certified. At order $7$,
$35$ of the $1{,}868$ instances are not. For EE at order $7$, every pair $(z,u)$ is certified except
the six pairs of $\mathrm{CP}(3)$ with an apex (\texttt{F]\char`~vw}), all of which have $|T|=4$.
Outputs: \texttt{out/cert\_5\_7.txt}, \texttt{out/cert\_ee7.txt}. So the reductions do not decide
the problem on their own. The cocktail party graphs are also the barrier examples for Fan-transfer
certificates in the descent programme (Lean: \texttt{cocktail\_fan\_barrier}).

\section{Formalisation}\label{sec:lean}

\texttt{lean/GallaiLean/L1Star.lean} (Lean~4, Mathlib, no \texttt{sorry}; \texttt{\#print axioms}
shows only \texttt{propext}, \texttt{Classical.choice}, \texttt{Quot.sound}; the full library
builds).
\begin{itemize}
\item \texttt{Admissible}, \texttt{AvoidDecomp}: admissible sets, good instances.
\texttt{L1StarGen}: Conjecture~\ref{conj:gen}. \texttt{L1Blocked}: every blocked instance is good.
\item \texttt{avoid\_of\_leafSep}: Lemma~\ref{lem:merge}. \texttt{avoid\_of\_apexSep}:
Lemma~\ref{lem:hang}(ii). Lemma~\ref{lem:hang}(i) is \texttt{PathDecomp.liftLeaf}~\cite{PE}.
\item \texttt{l1\_case\_adj}, \texttt{l1\_case\_target}, \texttt{l1\_case\_target\_adj},
\texttt{l1\_case\_big}: Theorem~\ref{thm:reductions}(a)--(d); \texttt{l1\_case\_leaf}:
Proposition~\ref{prop:allodd}. Both \texttt{l1\_case\_target} and \texttt{l1\_case\_target\_adj}
assume $c\notin N(g)$, as the proof of Theorem~\ref{thm:blocked} needs. Lemma~\ref{lem:protect} is
not a separate Lean lemma; it is inlined in \texttt{l1\_case\_target}.
\item \texttt{l1\_strong}, \texttt{l1starGen\_of\_blocked}: Theorem~\ref{thm:blocked}.
\texttt{l1star\_of\_gen}: Conjecture~\ref{conj:gen} implies \texttt{L1Star}.
\texttt{ee\_of\_L1StarGen}, \texttt{ee\_of\_l1Blocked}: Corollary~\ref{cor:eeblocked}.
\end{itemize}
Not formalised: Section~\ref{sec:more}, Proposition~\ref{prop:te3}, the examples.

\section{Open problems}\label{sec:open}
\begin{enumerate}
\item The blocked case of $L_1^*$, which by Corollary~\ref{cor:eeblocked} completes EE. By
Corollary~\ref{cor:mincex} the two shapes to handle are: $g\in A$ with every neighbour of $g$
adjacent to $A\setminus\{g\}$ (if $d(g)=2$, $N(g)\subseteq N(c)$ is an independent pair); and $A$
separating $c$ from $g$.
\item The two-even conjecture (Conjecture~\ref{conj:te3}), a statement of independent interest and a
consequence of the local form. Lov\'asz's Case~I step moves one even vertex at a time, and walks of
such steps do not always reach a smaller instance at distance three. At order $8$, one graph needs a
different argument (\texttt{G?Bvvo}, \texttt{te3walk.py}); so does one graph at order $6$
(\texttt{EEz\_}).
\item Is there a separation-type invariant, a property of the vertex sets of \emph{all} paths, that
implies $V(Q_c)\subseteq N[c]\cup N(g)$ and survives Fan's lemma? The one-sided form above is not
such an invariant, which is why $c$'s end must be protected at every step.
\end{enumerate}

\begin{thebibliography}{9}
\bibitem{Fan05} G.~Fan, Path decompositions and Gallai's conjecture, J.~Combin.\ Theory Ser.~B 93 (2005)
117--125.
\bibitem{L68} L.~Lov\'asz, On covering of graphs, in: Theory of Graphs (Proc.\ Colloq.\ Tihany, 1966),
Academic Press, New York, 1968, 231--236.
\bibitem{MS} J.~Williams, Separating many pairs in Lov\'asz decompositions: classes, the leaf merge, and
the Case~I step, manuscript (\texttt{manuscripts/multi-sep}), 2026.
\bibitem{PE} J.~Williams, Prescribed end edges in Lov\'asz decompositions with one even vertex,
manuscript (\texttt{manuscripts/prescribed-end}), 2026.
\bibitem{RD} J.~Williams, Route descent: transversal realisability, minimum $T$-joins, and Lov\'asz
decompositions that separate two vertices, manuscript (\texttt{manuscripts/route-descent}), 2026.
\end{thebibliography}
\end{document}
